% Appendix E -- The decomposition when efficiency lasts longer (draft, 29 October 2026)
%   tables    replication/output/tables/tab-payoff-horizon-varrho.tex
%             (code/4-exhibits/tab_payoff_horizon.R),
%             tab-flows-accounting-varrho0.96.tex (tab_flows_accounting.R -varrho0.96),
%             tab-mechanisms-varrho.tex (tab_mechanisms.R)
%   runs      replication/data/model/r0/lifecycle/*-varrho0.96.csv and
%             *-varrho0.96-disable_efficient_stock_purchase1.csv (simulate_lifecycle.py, CPU)
%   in-place  25-34, without A-C homes for sale: tab-flows-accounting-varrho0.96-
%             disable_efficient_stock_purchase1.csv (not a paper table)
%   parked    whether re-estimating with varrho fixed at 0.96 would change the
%             decomposition (author, 29 October 2026; to come back to)
\section{A Longer-Lasting Payoff}\label{app:varrho}

The retention of efficiency $\varrho$ is imprecisely estimated and lies close
to the lower end of its search range (Section~\ref{sec:estimates}). At the
estimate, $\varrho=0.84$, an efficient home keeps its rating for about six
years on average. This appendix repeats the three calculations of
Section~\ref{sec:results} at $\varrho=0.96$, at which a home keeps its rating
for 25 years. The other parameters stay at the estimates of
Table~\ref{tab:estimates}, and households re-optimize. The networks of
Appendix~\ref{app:computation} take $\varrho$ as an input and were trained
over its search range, $[0.80,0.99]$, so the model need not be solved again.
The calculations ask how the decomposition changes when efficiency lasts
longer, not whether the model still fits the data: at $\varrho=0.96$ the
efficient share is 43.1, 45.9, and 36.8 percent of owners at 25--34, 35--64,
and 65--75, about twice the shares in the data.

\subsection{The payoff horizon.}\label{app:varrho-horizon}
Table~\ref{tab:payoff-horizon-varrho} computes the horizon of
Section~\ref{sec:results-horizon} at both values of $\varrho$. Column~(1),
with discounting and survival only, does not depend on $\varrho$. Once the
loss of efficiency is added, the horizon at 75 is 83 percent of the horizon
at 25 at the estimate (column~(2)), but 55 percent at $\varrho=0.96$
(column~(5)): when efficiency lasts, survival starts to matter. Moving undoes
most of this. With the model's probabilities of moving, the horizon at
$\varrho=0.96$ is 6.2 years at 25 and 5.4 at 75 (column~(6)); with those in
Understanding Society it is 6.7 and 7.0 years (column~(7)). Most owners move
long before a rating that lasts 25 years wears off, so the horizon stays
short and close to flat in age.

\subsection{The accounting.}\label{app:varrho-accounting}
Table~\ref{tab:flows-accounting-varrho} repeats
Table~\ref{tab:flows-accounting} at $\varrho=0.96$. Buying an efficient home
is still the main way into the stock and the only one that falls steeply
with age: 8.8 per 100 owners at 25--34, 4.4 at 35--64, and 1.7 at 65--75,
while renovating in place adds 1.0 to 1.1 at every age (Panel~A). The loss of
efficiency is now a smaller outflow, 1.5 to 2.1 per 100 owners against 2.3 to
4.6 at the estimate, so older owners keep more of what they bought when
young, and the stock falls less with age than the inflow. In Panel~B
purchases fall from 15.0 per 100 owners at 25--34 to 3.8 at 65--75, close to
the 15.0 and 4.5 at the estimate. Unlike at the estimate, the share of those
purchases that are of efficient homes also falls with age, from 58 percent to
45 percent. Panel~C attributes the fall in the inflow. Purchases account for
90 percent of the fall from 25--34 and 89 percent of the fall from 35--64,
against 99 and 103 percent at the estimate. The share of efficient
purchases, the margin on which the horizon acts, accounts for 17 and 13
percent. Rebuilt from these margins, the efficient share at 65--75 would be
72.3 percent instead of 36.8 with purchases held at their rate at 25--34, and
41.6 percent with the share of efficient purchases held at its rate at
25--34.

\subsection{Closing the window.}\label{app:varrho-window}
Table~\ref{tab:mechanisms-varrho} removes efficient homes from the market at
both values of $\varrho$. A buyer can still buy an inefficient home and
renovate it at once, an owner can still renovate in place, and households
re-optimize. At $\varrho=0.96$ the efficient share falls to 13.8, 10.3, and
12.4 percent at 25--34, 35--64, and 65--75, and the gap between 25--34 and
65--75 shrinks from 6.2 points to 1.4 (column~(4)). At the estimate it shrinks
from 9.6 points to 1.3. Households that cannot buy efficiency renovate more:
at 25--34, renovation in place rises from 1.0 to 1.7 per 100 owners a year,
and renovation at purchase from 0.3 to 0.6. Yet the gap that remains without
the window is about the same at both values of $\varrho$. A longer-lasting
payoff adds little to the part of the gradient that does not run through
purchases.

The three calculations point the same way as at the estimate. A
longer-lasting payoff gives the horizon a role, through the share of buyers
who choose an efficient home, but the window still accounts for most of the
gradient.

\begin{table}[htbp]
  \centering
  \caption[The payoff horizon when efficiency lasts longer]{The
  payoff horizon when efficiency lasts longer}\label{tab:payoff-horizon-varrho}
  \begin{threeparttable}
    \small
    \pkgtable{tab-payoff-horizon-varrho}
    \begin{tablenotes}[flushleft]\footnotesize
      \item \textit{Notes:} The horizon $H_j$ of
      equation~\eqref{eq:payoff-horizon}, in years, as in
      Table~\ref{tab:payoff-horizon}. Columns~(1)--(4) are those of
      Table~\ref{tab:payoff-horizon}, at the estimates. Columns~(5)--(7) set
      $\varrho=0.96$ and keep the other parameters at the estimates; the
      model's probability that an owner of an A--C home moves, in
      column~(6), is from the simulation at $\varrho=0.96$. Column~(1) does
      not depend on $\varrho$. The last row is the horizon at 75 as a
      percentage of the horizon at 25.
    \end{tablenotes}
  \end{threeparttable}
\end{table}

\begin{table}[htbp]
  \centering
  \caption[Flows into and out of efficient ownership when efficiency lasts
  longer]{Flows into and out of efficient ownership when efficiency
  lasts longer}\label{tab:flows-accounting-varrho}
  \begin{threeparttable}
    \small
    \pkgtable{tab-flows-accounting-varrho0.96}
    \begin{tablenotes}[flushleft]\footnotesize
      \item \textit{Notes:} Simulated at $\varrho=0.96$, the other
      parameters at the estimates of Table~\ref{tab:estimates}, 65,536
      households, who re-optimize. Rows, panels, and the decomposition as in
      Table~\ref{tab:flows-accounting}.
    \end{tablenotes}
  \end{threeparttable}
\end{table}

\begin{table}[htbp]
  \centering
  \caption[The efficient share of owners without efficient homes for
  sale]{The efficient share of owners without efficient homes for
  sale}\label{tab:mechanisms-varrho}
  \begin{threeparttable}
    \small
    \pkgtable{tab-mechanisms-varrho}
    \begin{tablenotes}[flushleft]\footnotesize
      \item \textit{Notes:} Shares, gap, and ratio as in
      Table~\ref{tab:mechanisms-window}, other parameters at the estimates of
      Table~\ref{tab:estimates}. No A--C homes for sale: a buyer cannot buy
      an efficient home, but can still buy an inefficient home and renovate
      it at once, and an owner can still renovate in place; households
      re-optimize. EHS: the targets of Table~\ref{tab:estimates}.
      \item $^a$The rows at $\varrho=0.84$ are those of
      Table~\ref{tab:mechanisms-window}, simulated on a GPU; the rows at
      $\varrho=0.96$ are simulated on a CPU. At the estimates the CPU gives
      shares of 22.5, 17.7, and 12.7 percent; the two differ by simulation
      noise only.
    \end{tablenotes}
  \end{threeparttable}
\end{table}
